\section{Introduction}\label{sec:introduction}

Applications that serve users in several locations are increasingly deployed on more than one Kubernetes cluster. The clusters differ in capacity, in current load and in their network distance to the users. Orchestration across such clusters has been studied as part of the cloud-to-things continuum \cite{ullah2023orchestration}, and several systems coordinate independent clusters, among them the archived KubeFed project \cite{kubefed_repo}, Liqo \cite{iorio2023liquid} and Karmada \cite{karmada_repo,karmada_docs_architecture}. In Karmada, a central control plane receives a workload together with a propagation policy and decides on which member clusters the workload runs. This decision is made by the Karmada scheduler, which, unlike the Kubernetes scheduler, selects clusters rather than nodes.

The information used by this decision is limited. The Karmada scheduler evaluates API availability, cluster affinity, taints and an estimate of how many replicas each cluster can still accept. It does not use any measure of the network path between the clients of a workload and the candidate clusters. For latency-sensitive services, two clusters with equal free capacity can therefore receive the same preference although one of them is much farther from the clients. Tail latency matters in this setting, because response-time variability at the tail dominates the experience of interactive services at scale \cite{dean2013tail}.

Network-aware scheduling has been studied extensively at node level inside a single cluster \cite{santos2019networkaware,santos2023diktyo,wojciechowski2021netmarks,centofanti2023latency,ragavan2025nas}, and recent work has combined resource and latency information for selecting clusters \cite{tabatabaei2025cnf,bregar2026cosmospan,chiaro2024lais}. Combining the two kinds of information is therefore not new. What remains open, as discussed in \cref{sec:related}, is how such a combination fits into the mechanisms Karmada already provides. Existing proposals either delegate the decision to an external optimiser \cite{tabatabaei2025cnf} or replace Karmada's cluster scheduler with their own \cite{bregar2026cosmospan}. None of them specifies what happens when the network measurements are missing or out of date.

This paper addresses that question as a design study. It does not report an implementation or measurements; these are planned as a subsequent experimental study. The research questions are:
\begin{itemize}
  \item RQ1: How can resource and network information be incorporated into Karmada's cluster selection while keeping the existing scheduling behaviour as the foundation?
  \item RQ2: How should mandatory placement constraints be separated from optimisation objectives?
  \item RQ3: How can stale or unavailable network measurements be handled so that missing data is never treated as favourable?
  \item RQ4: Which trade-offs does the resulting design introduce?
\end{itemize}

The contributions of the paper are the following. First, it analyses the scheduling semantics of Karmada v1.19.1 at the level of the source code and identifies the properties that any custom cluster selection objective has to respect. Second, it proposes a design that keeps the hard eligibility conditions in Karmada's filter stage, reuses Karmada's own replica estimators for the capacity condition, and adds client-to-cluster round-trip time as a soft scoring signal with explicit validity and fallback rules. Third, it illustrates the decisions of the design analytically and proposes a protocol for validating it experimentally.

The paper is organised as follows. \Cref{sec:related} reviews related work. \Cref{sec:problem} formulates the problem. \Cref{sec:architecture} presents the architecture and \cref{sec:method} the cluster selection method. \Cref{sec:algorithm} describes the algorithm and its integration with Karmada. \Cref{sec:analysis} analyses the design, \cref{sec:validation} proposes the experimental validation, \cref{sec:limitations} states the limitations and \cref{sec:conclusions} concludes.
